\documentclass[11pt]{report}
\usepackage[margin=1in]{geometry}
\usepackage{amsmath}
\usepackage{amssymb}
\usepackage{amsthm}
\usepackage{graphicx}
\usepackage{hyperref}
\usepackage{comment}

\newtheorem{theorem}{Theorem}
\newtheorem{lemma}{Lemma}
\newtheorem{proposition}{Proposition}
\newtheorem{corollary}{Corollary}
\newtheorem{assumption}{Assumption}
\usepackage{tikz-cd}
\usepackage{multicol}

\setlength{\parindent}{0pt}
\setlength{\parskip}{1\baselineskip}

\begin{document}

\chapter{Why Simple Averaging Works for Transformer Gradient Structure}

\section{Gradients of Linear Layers in Transformers}
\label{sec:gradients_linear_layers}

We begin by formalizing the relationship between model parameters, activations, and
gradients for a linear projection in a transformer layer.
This relationship is central to understanding how low-rank and spiky structure
emerges in gradient matrices and how it influences optimization dynamics.

\paragraph{Linear projection and model parameters.}
Consider a single transformer layer and one of its linear projections, for example
the key projection matrix
\[
W_k \in \mathbb{R}^{d_{\text{out}} \times d_{\text{in}}}.
\]
In standard transformer architectures, $d_{\text{out}} = d_{\text{in}} = d$.
Given an input hidden state
\[
x \in \mathbb{R}^{d_{\text{in}}},
\]
the key vector is computed as
\[
k = W_k x \in \mathbb{R}^{d_{\text{out}}}.
\]

The rows of $W_k$ correspond to output features of the key vector, while the columns
correspond to input features of the hidden state.
The parameter matrix $W_k$ therefore defines how input features are linearly combined
to produce output features.

\paragraph{Token-level activations.}
During training, the model processes a batch of tokens.
For token $j$, let
\[
x^{(j)} \in \mathbb{R}^{d_{\text{in}}}
\]
denote the input activation to the key projection.
Stacking activations across a batch of $n$ tokens yields the activation matrix
\[
X \in \mathbb{R}^{n \times d_{\text{in}}},
\qquad
X_{j,:} = \bigl(x^{(j)}\bigr)^\top.
\]
The activation matrix encodes the representation geometry of the model at the input
of the linear layer.

\paragraph{Backpropagated error signals.}
Let
\[
\delta_k^{(j)} = \frac{\partial \mathcal{L}}{\partial k^{(j)}}
\in \mathbb{R}^{d_{\text{out}}}
\]
denote the backpropagated error signal associated with token $j$ at the key
projection.
Stacking these error vectors gives
\[
\Delta_k \in \mathbb{R}^{n \times d_{\text{out}}},
\qquad
(\Delta_k)_{j,:} = \bigl(\delta_k^{(j)}\bigr)^\top.
\]

\paragraph{Gradient matrix.}
During backpropagation, the gradient of the loss $\mathcal{L}$ with respect to $W_k$
has the same shape as the parameter matrix:
\[
G \;\triangleq\; \nabla_{W_k} \mathcal{L}
\;\in\;
\mathbb{R}^{d_{\text{out}} \times d_{\text{in}}}.
\]
For a single token $j$, the contribution to the gradient has the outer-product form
\[
\nabla_{W_k} \mathcal{L}^{(j)}
=
\delta_k^{(j)} \bigl(x^{(j)}\bigr)^\top
\;\in\;
\mathbb{R}^{d_{\text{out}} \times d_{\text{in}}}.
\]
Summing over all $n$ tokens in the batch, the full gradient matrix can be written as
\[
G
=
\sum_{j=1}^{n}
\delta_k^{(j)} \bigl(x^{(j)}\bigr)^\top
=
\Delta_k^\top X.
\]
Thus, although $G$ is a single $d_{\text{out}} \times d_{\text{in}}$ matrix, it
implicitly aggregates correlations between activation directions and error
directions across many tokens and contexts.

\paragraph{Connection to stochastic gradient descent.}
Under vanilla stochastic gradient descent (SGD), the parameter update at training
step $t$ is
\[
W_k^{(t+1)}
=
W_k^{(t)} - \eta \, G^{(t)},
\]
where $\eta > 0$ is the learning rate.
Substituting the gradient expression yields
\[
W_k^{(t+1)}
=
W_k^{(t)}
-
\eta
\sum_{j=1}^{n}
\delta_k^{(j,t)} \bigl(x^{(j,t)}\bigr)^\top.
\]

Unrolling SGD over multiple steps shows that the parameter matrix can be interpreted
as a time-integrated sum of gradient outer products:
\[
W_k^{(T)}
=
W_k^{(0)}
-
\eta
\sum_{t=1}^{T}
\sum_{j}
\delta_k^{(j,t)} \bigl(x^{(j,t)}\bigr)^\top.
\]
Consequently, the structure of the parameter matrix reflects the accumulated
correlations between activation patterns and backpropagated error signals over the
course of training.

\paragraph{Summary.}
Activations $X$ describe what features are present in the data, gradients $G$ describe
how the model parameters are being pushed to change in response to those features,
and the parameter matrix $W_k$ encodes the historical integration of these gradient
signals.
In the following sections, we focus on the gradient matrix $G$ as the primary object
for analyzing optimization geometry and for designing simple, GPU-friendly
approximations to its dominant low-dimensional structure.
\section{Assumptions on Grammar and Semantic Correlation}
\label{sec:foundational_assumptions}

The averaging-based approximation developed in this chapter relies on a set of
assumptions about how grammatical and semantic information is distributed across
tokens in training data.
These assumptions are not specific to a particular model architecture, but rather
reflect generic statistical properties of natural language corpora and modern
training pipelines.
They are critical for ensuring that the proposed method isolates shared linguistic
structure while preserving the semantic correlations that must be learned.

\paragraph{Token indexing and notation.}
We index tokens by a pair $(j,t)$, where $j$ denotes the data sequence (e.g., a
sentence or document) and $t$ denotes the token position within that sequence.
Gradients, activations, and error signals are aggregated over tokens in a training
batch, which typically contains multiple sequences.

\paragraph{Assumption 1: Shared grammatical structure.}
We assume that all tokens, regardless of sequence or position, share a common set
of grammatical and structural regularities.
These include syntactic patterns, function words, positional conventions, and other
linguistic features that recur ubiquitously across natural language.
As a result, grammatical components of activations and gradients are \emph{globally
aligned} across tokens and sequences, and induce coherent, low-dimensional structure
in aggregated gradient matrices.

\paragraph{Assumption 2: Intra-sequence semantic correlation.}
Tokens within the same data sequence exhibit strong semantic correlation.
They share entities, topics, discourse context, and compositional meaning.
Consequently, semantic components of activations and gradients are correlated across
tokens belonging to the same sequence.
These correlations are essential for learning meaningful representations and must
not be suppressed by any approximation used during training.

\paragraph{Assumption 3: Inter-sequence semantic independence under data shuffling.}
We assume that the training data is well shuffled, so that different sequences in a
batch are drawn independently.
Under this assumption, semantic components of tokens belonging to \emph{different}
sequences have very low correlation in expectation.
While individual tokens may have large semantic gradients, their directions vary
across sequences and therefore cancel when aggregated across a sufficiently diverse
batch.

\paragraph{Extension to long-context training.}
In long-context settings, a single sequence may contain thousands of tokens.
Although tokens that are close in position are strongly correlated semantically,
semantic correlation decays with distance.
We therefore assume that tokens within the same sequence but separated by a
sufficiently large context window can be treated as approximately semantically
independent.
In this regime, averaging across such distant tokens behaves similarly to averaging
across independent sequences, while still preserving local semantic structure.

\paragraph{Implications for averaging-based approximation.}
Taken together, these assumptions imply a clear separation of roles:
\begin{itemize}
    \item Grammar-induced components are shared across tokens and sequences and
    reinforce coherently under aggregation, producing spiky, low-dimensional
    structure in gradient matrices.
    \item Semantic components are correlated within a sequence but largely
    uncorrelated across sequences (or across sufficiently distant positions in
    long-context settings).
    \item Averaging across tokens drawn from different sequences (or distant
    positions) suppresses independent semantic variation while preserving shared
    grammatical structure.
\end{itemize}

This separation is essential for the proposed approximation method.
It allows simple averaging operations to approximate the dominant, grammar-driven
low-rank structure of the agreement matrix, while leaving intra-sequence semantic
correlations intact.
As a result, the method isolates the spiky structure that dominates optimization
dynamics without interfering with the learning of semantic features.

\paragraph{Scope.}
Throughout this chapter, all averaging operations are implicitly understood to be
performed over tokens that are approximately semantically independent under the
assumptions above.
Violations of these assumptions---for example, highly correlated batches or
short-context averaging---may reduce the effectiveness of the approximation and are
outside the scope of the present analysis.
\newpage

\section{Empirical Validation of Foundational Assumptions}
\label{sec:empirical_validation_assumptions}

The analysis in this chapter is built upon a set of foundational assumptions
regarding the distribution of grammatical and semantic information across tokens,
sequences, and training batches.
While these assumptions are motivated by linguistic structure and common training
practice, their validity is ultimately an empirical question.
The purpose of the experiments in this section is therefore to quantitatively
evaluate these assumptions in realistic transformer models and datasets.

Specifically, we aim to measure (i) the degree to which grammatical structure induces
globally shared, low-dimensional patterns in activations and gradients; (ii) the
strength of semantic correlation among tokens within the same sequence; and (iii)
the extent to which semantic correlations diminish across independently shuffled
sequences or sufficiently distant positions in long-context settings.
By explicitly characterizing these correlation patterns, we seek to verify that
simple averaging across appropriately chosen tokens suppresses independent semantic
variation while preserving the shared grammatical structure that dominates gradient
geometry.
This empirical validation is critical for establishing when and why the proposed
averaging-based approximation accurately captures the spiky, low-rank structure
driving optimization dynamics.

\paragraph{Preliminary spectral alignment observation.}
Given an activation tensor $a\in\mathbb{R}^{\texttt{B}\times\texttt{L}\times\texttt{H}}$, we reshape it into
$b\in\mathbb{R}^{(\texttt{B}\cdot\texttt{L})\times\texttt{H}}$ and form a ``column-mean matrix''
$\mathrm{ColMean}(b)$ by averaging each column of $b$ and broadcasting the resulting mean vector to match the shape of $b$.
We then compute the SVD of $b$ and measure the (absolute) cosine similarity between $\mathrm{ColMean}(b)$ and the reconstruction induced by the $k$-th singular component.
Empirically, we observe that the leading $\approx 1.5\%$ of the spectral subspace exhibits high alignment with $\mathrm{ColMean}(b)$,
whereas the remaining $\approx 98.5\%$ is nearly orthogonal.
Moreover, the alignment strength depends on how the column mean is estimated: averaging over all sentences in the batch yields the
strongest alignment, followed by averaging over randomly sampled tokens across sentences, while using a single sentence produces the
weakest alignment.

\paragraph{Spectral alignment under different mean estimators.}
We formalize and extend the above observation by comparing three estimators of the column mean.
Let $a\in\mathbb{R}^{B\times L\times H}$ denote the activation tensor, and let $b\in\mathbb{R}^{(BL)\times H}$ be the reshaped matrix obtained by
flattening the batch and sequence dimensions.
Given a mean vector $\mu\in\mathbb{R}^{H}$, we define the corresponding broadcasted ``column-mean matrix'' as
$\mathrm{ColMean}_{\mu}(b)=\mathbf{1}\mu^{\top}\in\mathbb{R}^{(BL)\times H}$.
For each $k$, we compare $\mathrm{ColMean}_{\mu}(b)$ with the rank-$1$ reconstruction induced by the $k$-th singular component of $b$, i.e.,
$\sigma_k u_k v_k^{\top}$, using the absolute cosine similarity in Frobenius inner product.

We consider the following three choices of $\mu$:
(i) \emph{Batch mean} (Method~1): average each column of $b$ over all $BL$ tokens in the full batch;
(ii) \emph{Single-sequence mean} (Method~2): compute $\mu$ from one sequence only (e.g., averaging columns of $a[0,:,:]\in\mathbb{R}^{L\times H}$) and broadcast it to shape $(BL)\times H$;
(iii) \emph{Random token sampling across sequences} (Method~3): randomly sample a subset of tokens from each of the $B$ sequences to form a
set of token vectors, compute the column mean over this set, and broadcast.

Across all three estimators, we find a consistent spectral dichotomy: the leading $\approx 1.5\%$ of singular components exhibits high
alignment with $\mathrm{ColMean}_{\mu}(b)$, while the remaining $\approx 98.5\%$ is nearly orthogonal.
Moreover, the alignment strength follows a clear ordering,
\emph{batch mean} (Method~1) $>$ \emph{random token sampling} (Method~3) $>$ \emph{single-sequence mean} (Method~2),
which is consistent with the chapter's core hypothesis that the shared, low-rank structure is most reliably recovered when averaging over
many semantically diverse tokens while preserving the globally shared grammatical signal.
Figures~\ref{fig:colmean_alignment_method1}--\ref{fig:colmean_alignment_method3} report the cosine similarity curves for the three estimators.

\begin{figure}[t]
  \centering
  \includegraphics[width=0.92\linewidth]{figures/col_batch.png}
  \caption{Absolute cosine similarity between $\mathrm{ColMean}_{\mu}(b)$ (with $\mu$ estimated using the full batch) and the rank-$1$ matrix
  $\sigma_k u_k v_k^{\top}$ induced by the $k$-th singular component of $b$.}
  \label{fig:colmean_alignment_method1}
\end{figure}

\begin{figure}[t]
  \centering
  \includegraphics[width=0.92\linewidth]{figures/col_single.png}
  \caption{Same as Figure~\ref{fig:colmean_alignment_method1}, but with $\mu$ estimated from a single sequence (e.g., $a[0,:,:]$) and then
  broadcast to shape $(BL)\times H$.}
  \label{fig:colmean_alignment_method2}
\end{figure}

\begin{figure}[t]
  \centering
  \includegraphics[width=0.92\linewidth]{figures/col_sampled.png}
  \caption{Same as Figure~\ref{fig:colmean_alignment_method1}, but with $\mu$ estimated from randomly sampled tokens across sequences.}
  \label{fig:colmean_alignment_method3}
\end{figure}

\newpage

\section{What Structure Are We Approximating?}
\label{sec:target_structure}

Before introducing specific approximation methods, it is essential to clarify
\emph{what kind of structure} we aim to extract from transformer gradient matrices.
Although the gradients studied in this chapter often exhibit strong low-rank or
spiky spectra, our objective is \emph{not} to approximate an arbitrary low-rank
matrix in the classical numerical linear algebra sense.
Instead, we focus on a highly specific and structured form of low-dimensionality
that arises from shared linguistic regularities.

\paragraph{Agreement-driven structure.}
As established in Sections~\ref{sec:gradients_linear_layers} and
\ref{sec:foundational_assumptions}, the gradient matrix of a transformer linear
layer can be written as a sum of per-token outer products,
\[
G = \sum_{j,t} \delta^{(j,t)} \bigl(x^{(j,t)}\bigr)^\top.
\]
Under the assumptions that (i) grammatical structure is shared across tokens and
sequences, and (ii) semantic variation is largely independent across sequences,
the aggregated gradient naturally decomposes as
\[
G = A + R,
\]
where:
\begin{itemize}
    \item $A$ is a low-rank \emph{agreement component} induced by grammar-driven,
    globally aligned update directions; and
    \item $R$ is a high-rank residual capturing semantic and context-specific
    variation that cancels under appropriate aggregation.
\end{itemize}

The term \emph{agreement} emphasizes that $A$ arises from consistent, repeated
update signals shared across many tokens and sequences, rather than from generic
low-rank structure.

\paragraph{Why this is not generic low-rank approximation.}
In general, a low-rank matrix may have singular vectors that are sparse,
oscillatory, or orthogonal to simple probe directions.
Recovering such structure typically requires explicit singular value
decomposition or iterative orthogonalization.
In contrast, the agreement component $A$ considered here has several special
properties:
\begin{itemize}
    \item its singular vectors are dense and globally aligned,
    \item it induces non-vanishing row and column marginals,
    \item it survives simple aggregation operations such as summation or averaging.
\end{itemize}
These properties are a direct consequence of shared grammatical structure and
sequence-level semantic independence, not of low rank alone.

\paragraph{Implications for approximation.}
Because the agreement component $A$ produces coherent, non-canceling marginals,
it can be detected using extremely simple probes, such as multiplication by the
all-ones vector or by a small number of random vectors.
These probes are sufficient to recover the dominant subspace of $A$ up to scale,
without requiring explicit singular value decomposition or QR factorization.

Crucially, this approximation strategy is \emph{structure-exploiting} rather than
\emph{structure-agnostic}.
It is designed to isolate grammar-driven agreement spikes while leaving
intra-sequence semantic correlations intact, thereby preserving the information
necessary for learning semantic features.

\paragraph{Scope and limitations.}
Throughout this chapter, all approximation results should be interpreted in the
context of this agreement-driven structure.
The proposed averaging-based methods are not intended to recover arbitrary
low-rank matrices, nor to replace general-purpose SVD algorithms.
Rather, they provide a lightweight, GPU-friendly mechanism for approximating the
dominant optimization-driving structure that arises under realistic assumptions
about linguistic data and training pipelines.

With this target structure in mind, we now turn to averaging operations as simple
yet effective probes for isolating agreement-driven components in gradient
matrices.
\section{Averaging as a Probe for Agreement-Driven Structure}
\label{sec:averaging_probe}

Having clarified that our objective is to approximate grammar-driven agreement
structure rather than arbitrary low-rank matrices
(Section~\ref{sec:target_structure}),
we now introduce averaging operations as simple, GPU-friendly probes for isolating
this structure in gradient matrices.

\subsection{Physical Interpretation of Gradient Entries}
\label{sec:physical_gradient}

We briefly recall the physical meaning of the gradient matrix entries in order to
motivate the averaging operations that follow.

Consider the key projection in a transformer layer,
\[
k = W_k x,
\qquad
W_k \in \mathbb{R}^{d\times d},
\]
where $x \in \mathbb{R}^{d}$ is the input hidden state and $k \in \mathbb{R}^{d}$ is
the resulting key vector.
The gradient of the loss with respect to $W_k$,
\[
G = \nabla_{W_k} \mathcal{L} \in \mathbb{R}^{d\times d},
\]
has entries
\[
G_{i\ell} = \frac{\partial \mathcal{L}}{\partial W_{k,i\ell}},
\]
which measure how strongly the training data pushes the connection from input
feature $\ell$ to output feature $i$.

For a batch of tokens, the gradient aggregates per-token outer products,
\[
G = \sum_{j,t} \delta^{(j,t)} \bigl(x^{(j,t)}\bigr)^\top,
\]
where $\delta^{(j,t)}$ is the backpropagated error signal at the projection for
token $(j,t)$.
Thus, $G$ encodes correlations between activation directions and error directions,
aggregated across many tokens and contexts.

\subsection{Column Averaging: Primary Probe of Agreement}
\label{sec:column_averaging}

We now introduce column averaging as the primary operation for probing
agreement-driven structure.

Let $\mathbf{1} \in \mathbb{R}^{d}$ denote the vector of all ones, and define the
column-averaged vector
\[
u \;\triangleq\; G \mathbf{1} \in \mathbb{R}^{d},
\qquad
u_i = \sum_{\ell=1}^{d} G_{i\ell}.
\]

Substituting the token-level gradient expression yields
\[
u
=
\sum_{j,t}
\delta^{(j,t)} \bigl( (x^{(j,t)})^\top \mathbf{1} \bigr).
\]
Importantly, this expression remains a \emph{sum over tokens}.
The multiplication by $\mathbf{1}$ collapses feature-specific input variation into
a scalar weight for each token, while the aggregation over tokens suppresses
semantic variation under the assumptions of
Section~\ref{sec:foundational_assumptions}.

\paragraph{Interpretation.}
Column averaging answers the question:
\begin{quote}
\emph{Which output directions of this layer are being consistently pushed across
many tokens and sequences, regardless of which specific input features caused the
pressure?}
\end{quote}

Because grammatical structure is shared across sequences, grammar-induced gradient
components reinforce coherently in $\nu$.
In contrast, semantic components vary across sequences and cancel in expectation
when tokens are drawn from semantically independent sequences.
As a result, $\nu$ serves as a direct probe of the dominant agreement-driven update
direction.

\paragraph{Why column averaging works.}
The effectiveness of column averaging does not rely on low rank alone.
It relies on the fact that the agreement spike is both \emph{low rank} and \emph{coherent}:
its singular vectors are dense, and its row/column marginals (e.g., $A\mathbf{1}$ and
$A^\top\mathbf{1}$) do not vanish due to cancellation.
In contrast, the residual semantic component is high-rank and comparatively \emph{incoherent}
under batching: token-specific semantic features vary across sequences and cancel in
these first-order aggregated statistics.
This asymmetry is specific to grammar-driven agreement structure and does not hold
for arbitrary low-rank matrices.

\subsection{Row Averaging: Auxiliary Construction}
\label{sec:row_averaging}

For completeness, we define the row-averaged vector
\[
v \;\triangleq\; G^\top \mathbf{1} \in \mathbb{R}^{d},
\qquad
v_\ell = \sum_{i=1}^{d} G_{i\ell}.
\]
In token-level form,
\[
v
=
\sum_{j,t}
x^{(j,t)} \bigl( (\delta^{(j,t)})^\top \mathbf{1} \bigr).
\]

Row averaging aggregates gradient contributions associated with each input feature
across output dimensions.
Empirically and conceptually, this signal is typically weaker and noisier than the
column-averaged signal, as input features mix semantic and contextual information
more heavily.

We emphasize that row averaging is \emph{not} introduced as an independent probe of
agreement structure.
Rather, it plays an auxiliary role in completing a rank-one matrix reconstruction
once the dominant agreement direction has been identified via column averaging.

\subsection{Rank-One Reconstruction from Averages}
\label{sec:rank_one_from_averages}

Given the column-averaged vector $\nu$ and the row-averaged vector $v$, a natural
rank-one approximation of the gradient matrix is
\[
\widehat{G}_1
\;\triangleq\;
\frac{u v^\top}{\mathbf{1}^\top G \mathbf{1}}.
\]

This construction is algebraic rather than structural: it combines the dominant
left- and right-marginal signals into a rank-one matrix.
Its validity rests entirely on the agreement-driven structure described in
Section~\ref{sec:target_structure} and the cancellation assumptions of
Section~\ref{sec:foundational_assumptions}.

In particular, $\widehat{G}_1$ should be interpreted as an approximation to the
agreement component $A$, not as a generic low-rank approximation to $G$.
When these assumptions fail, the reconstruction may be inaccurate, even if $G$ is
low rank.

\paragraph{Summary.}
Averaging serves as a lightweight, GPU-native probe for isolating agreement-driven
structure in gradient matrices.
Column averaging is the primary mechanism by which shared grammatical update
directions are revealed, while row averaging plays a secondary, algebraic role in
constructing a rank-one approximation.
This asymmetric treatment reflects the underlying linguistic and statistical
structure of the data and is essential for the correctness of the proposed method.
%\section{Rank-One Agreement Approximation: What Works and Why (with Detailed Proofs)}
\label{sec:rank_one_agreement}

This section proves two statements that are both critical for the proposed approximation method:

\begin{enumerate}
    \item \textbf{Agreement capture.} The averaging-based rank-one estimator
    $\widehat{G}_1$ recovers the dominant grammar-driven \emph{agreement spike}
    $A_{\mathrm{agr}}$ up to a controlled error.
    \item \textbf{Semantic non-capture (yet learnability).} The same averaging probe
    does \emph{not} extract intra-sequence semantic structure: semantic contributions can be
    large at the token level and remain available to SGD, while still canceling in the
    specific aggregated statistics used by the probe.
\end{enumerate}

Throughout, we make the dependence on the foundational assumptions in
Section~\ref{sec:foundational_assumptions} explicit.

\subsection{Setup: agreement component plus semantic residual}
\label{sec:rank_one_setup}

Let $G\in\mathbb{R}^{d\times d}$ be the gradient matrix of a transformer linear layer.
We decompose
\begin{equation}
\label{eq:rank_one_decomp}
G \;=\; A_{\mathrm{agr}} \;+\; R_{\mathrm{sem}},
\end{equation}
where $A_{\mathrm{agr}}$ is the grammar-driven \emph{agreement component} and $R_{\mathrm{sem}}$
is the semantic/context-specific residual.

\paragraph{Rank-one agreement model.}
Assume the agreement component is rank one:
\begin{equation}
\label{eq:rank_one_A}
A_{\mathrm{agr}} \;=\; \sigma a b^\top,
\qquad
a,b\in\mathbb{R}^d\setminus\{0\},
\qquad
\sigma>0.
\end{equation}

\paragraph{Visibility under the averaging probe.}
Let $\mathbf{1}\in\mathbb{R}^d$ be the all-ones vector and define
\begin{equation}
\label{eq:alpha_beta}
\alpha \triangleq \mathbf{1}^\top a,
\qquad
\beta \triangleq \mathbf{1}^\top b.
\end{equation}
Assume $\alpha\neq 0$ and $\beta\neq 0$, so that $A_{\mathrm{agr}}$ is visible to this probe.

\subsection{Semantic cancellation under aggregation: formal sequence-level proof}
\label{sec:semantic_cancellation_proof}

We now prove a concrete, sequence-level cancellation statement for the probe statistics
$R_{\mathrm{sem}}\mathbf{1}$, $R_{\mathrm{sem}}^\top\mathbf{1}$, and $\mathbf{1}^\top R_{\mathrm{sem}}\mathbf{1}$.

\subsubsection{Batch model and assumptions}

Consider a batch consisting of $S$ sequences. For sequence $s\in\{1,\dots,S\}$,
let $\mathcal{T}_s$ be the set of token indices included from that sequence
(e.g.\ all tokens, or a subset of sufficiently distant positions in long-context settings).
Write the batch gradient as
\begin{equation}
\label{eq:G_seq_sum}
G \;=\; \sum_{s=1}^S \sum_{t\in\mathcal{T}_s} \delta^{(s,t)} \bigl(x^{(s,t)}\bigr)^\top.
\end{equation}
Conceptually decompose each token-level contribution as
\begin{equation}
\label{eq:token_decomp}
\delta^{(s,t)} \bigl(x^{(s,t)}\bigr)^\top
\;=\;
A_{\mathrm{agr}}^{(s,t)} \;+\; R_{\mathrm{sem}}^{(s,t)}.
\end{equation}
Define the sequence-level semantic residual sum
\begin{equation}
\label{eq:Zs_def}
Z_s \;\triangleq\; \sum_{t\in\mathcal{T}_s} R_{\mathrm{sem}}^{(s,t)},
\qquad\text{so that}\qquad
R_{\mathrm{sem}} \;=\; \sum_{s=1}^S Z_s.
\end{equation}

We now formalize the statistical content of Section~\ref{sec:foundational_assumptions}:

\begin{assumption}[Sequence-level semantic independence and zero-mean]
\label{assump:Zs}
Conditioned on model parameters and the batch construction rule, the random matrices
$\{Z_s\}_{s=1}^S$ satisfy:
\begin{enumerate}
\item (\emph{Inter-sequence independence}) $Z_1,\dots,Z_S$ are independent.
\item (\emph{Zero-mean across sequences}) $\mathbb{E}[Z_s]=0$ for each $s$.
\end{enumerate}
For long-context training, the same assumption applies when $\mathcal{T}_s$ is chosen to
include only sufficiently distant positions so that semantic dependence within $\mathcal{T}_s$
is weak, i.e.\ $Z_s$ aggregates approximately independent semantic blocks.
\end{assumption}

\begin{assumption}[Finite second moments of probe responses]
\label{assump:moments}
There exist constants $\sigma_u^2,\sigma_v^2,\sigma_s^2<\infty$ such that for all $s$,
\begin{equation}
\label{eq:moment_bounds}
\mathbb{E}\|Z_s\mathbf{1}\|_2^2 \le \sigma_u^2,
\qquad
\mathbb{E}\|Z_s^\top\mathbf{1}\|_2^2 \le \sigma_v^2,
\qquad
\mathbb{E}\bigl(\mathbf{1}^\top Z_s \mathbf{1}\bigr)^2 \le \sigma_s^2.
\end{equation}
\end{assumption}

\subsubsection{Lemma: cancellation of semantic residual in the probe statistics}

\begin{lemma}[Semantic cancellation under aggregation (detailed)]
\label{lem:semantic_cancellation}
Let $R_{\mathrm{sem}}=\sum_{s=1}^S Z_s$ with $\{Z_s\}$ satisfying
Assumptions~\ref{assump:Zs}--\ref{assump:moments}. Then:
\begin{align}
\label{eq:cancel_mean}
\mathbb{E}[R_{\mathrm{sem}}\mathbf{1}] &= 0,
\qquad
\mathbb{E}[R_{\mathrm{sem}}^\top\mathbf{1}] = 0,
\qquad
\mathbb{E}[\mathbf{1}^\top R_{\mathrm{sem}}\mathbf{1}] = 0,\\
\label{eq:cancel_second_moment}
\mathbb{E}\|R_{\mathrm{sem}}\mathbf{1}\|_2^2 &\le S\,\sigma_u^2,
\qquad
\mathbb{E}\|R_{\mathrm{sem}}^\top\mathbf{1}\|_2^2 \le S\,\sigma_v^2,
\qquad
\mathbb{E}\bigl(\mathbf{1}^\top R_{\mathrm{sem}}\mathbf{1}\bigr)^2 \le S\,\sigma_s^2.
\end{align}
Consequently, for any $\delta\in(0,1)$, with probability at least $1-\delta$,
\begin{equation}
\label{eq:cancel_hp}
\|R_{\mathrm{sem}}\mathbf{1}\|_2 \le \sqrt{\frac{S}{\delta}}\,\sigma_u,\quad
\|R_{\mathrm{sem}}^\top\mathbf{1}\|_2 \le \sqrt{\frac{S}{\delta}}\,\sigma_v,\quad
|\mathbf{1}^\top R_{\mathrm{sem}}\mathbf{1}| \le \sqrt{\frac{S}{\delta}}\,\sigma_s.
\end{equation}
\end{lemma}

\begin{proof}
We prove each claim explicitly.

\paragraph{Mean identities.}
Using \eqref{eq:Zs_def} and linearity of expectation,
\[
\mathbb{E}[R_{\mathrm{sem}}\mathbf{1}]
=
\mathbb{E}\Big[\sum_{s=1}^S Z_s\mathbf{1}\Big]
=
\sum_{s=1}^S \mathbb{E}[Z_s]\mathbf{1}
=
0,
\]
where the last equality uses $\mathbb{E}[Z_s]=0$.
The proofs of $\mathbb{E}[R_{\mathrm{sem}}^\top\mathbf{1}]=0$ and
$\mathbb{E}[\mathbf{1}^\top R_{\mathrm{sem}}\mathbf{1}]=0$ are identical.

\paragraph{Second-moment bound for $R_{\mathrm{sem}}\mathbf{1}$.}
Let $U_s \triangleq Z_s\mathbf{1}\in\mathbb{R}^d$. Then $R_{\mathrm{sem}}\mathbf{1}=\sum_{s=1}^S U_s$.
Expand the squared norm:
\[
\Big\|\sum_{s=1}^S U_s\Big\|_2^2
=
\sum_{s=1}^S \|U_s\|_2^2
+
2\sum_{1\le s<t\le S} U_s^\top U_t.
\]
Taking expectations gives
\[
\mathbb{E}\Big\|\sum_{s=1}^S U_s\Big\|_2^2
=
\sum_{s=1}^S \mathbb{E}\|U_s\|_2^2
+
2\sum_{s<t} \mathbb{E}[U_s^\top U_t].
\]
By independence of $Z_s$ and $Z_t$, $U_s$ and $U_t$ are independent; since $\mathbb{E}[U_s]=0$,
\[
\mathbb{E}[U_s^\top U_t]
=
\mathbb{E}[U_s]^\top \mathbb{E}[U_t]
=
0.
\]
Thus,
\[
\mathbb{E}\|R_{\mathrm{sem}}\mathbf{1}\|_2^2
=
\sum_{s=1}^S \mathbb{E}\|U_s\|_2^2
=
\sum_{s=1}^S \mathbb{E}\|Z_s\mathbf{1}\|_2^2
\le
S\,\sigma_u^2,
\]
using \eqref{eq:moment_bounds}. This proves the first inequality in \eqref{eq:cancel_second_moment}.
The other two second-moment inequalities follow by the same argument with
$V_s \triangleq Z_s^\top\mathbf{1}$ and $W_s \triangleq \mathbf{1}^\top Z_s\mathbf{1}$.

\paragraph{High-probability bounds.}
Apply Markov's inequality to nonnegative random variables:
\[
\mathbb{P}\!\left(\|R_{\mathrm{sem}}\mathbf{1}\|_2 \ge \tau\right)
=
\mathbb{P}\!\left(\|R_{\mathrm{sem}}\mathbf{1}\|_2^2 \ge \tau^2\right)
\le
\frac{\mathbb{E}\|R_{\mathrm{sem}}\mathbf{1}\|_2^2}{\tau^2}
\le
\frac{S\sigma_u^2}{\tau^2}.
\]
Setting $\tau=\sqrt{S\sigma_u^2/\delta}$ yields
$\mathbb{P}(\|R_{\mathrm{sem}}\mathbf{1}\|_2 \ge \sqrt{S/\delta}\,\sigma_u)\le\delta$.
The other two bounds follow identically.
\end{proof}

\subsubsection{From cancellation to the $\varepsilon$-style inequalities}

To relate Lemma~\ref{lem:semantic_cancellation} to the relative inequalities used by the estimator,
we also need that the agreement component reinforces \emph{linearly} with the number of independent
sequences, reflecting shared grammar structure.

\begin{assumption}[Agreement reinforcement in probe marginals]
\label{assump:agreement_reinforce}
There exist constants $m_u,m_v,m_s>0$ such that
\begin{equation}
\label{eq:agreement_linear_scaling}
\|A_{\mathrm{agr}}\mathbf{1}\|_2 \ge S\,m_u,\qquad
\|A_{\mathrm{agr}}^\top\mathbf{1}\|_2 \ge S\,m_v,\qquad
|\mathbf{1}^\top A_{\mathrm{agr}}\mathbf{1}| \ge S\,m_s.
\end{equation}
\end{assumption}

\begin{corollary}[Relative semantic cancellation for the probe statistics]
\label{cor:eps_condition}
Under Lemma~\ref{lem:semantic_cancellation} and Assumption~\ref{assump:agreement_reinforce},
for any $\delta\in(0,1)$, with probability at least $1-\delta$,
\begin{equation}
\label{eq:residual_cancel}
\|R_{\mathrm{sem}}\mathbf{1}\|
\le
\varepsilon_u \|A_{\mathrm{agr}}\mathbf{1}\|,
\qquad
\|R_{\mathrm{sem}}^\top\mathbf{1}\|
\le
\varepsilon_v \|A_{\mathrm{agr}}^\top\mathbf{1}\|,
\qquad
|\mathbf{1}^\top R_{\mathrm{sem}}\mathbf{1}|
\le
\varepsilon_s |\mathbf{1}^\top A_{\mathrm{agr}}\mathbf{1}|,
\end{equation}
where
\[
\varepsilon_u \triangleq \frac{1}{\sqrt{S}}\cdot \frac{\sigma_u}{m_u}\cdot\frac{1}{\sqrt{\delta}},
\qquad
\varepsilon_v \triangleq \frac{1}{\sqrt{S}}\cdot \frac{\sigma_v}{m_v}\cdot\frac{1}{\sqrt{\delta}},
\qquad
\varepsilon_s \triangleq \frac{1}{\sqrt{S}}\cdot \frac{\sigma_s}{m_s}\cdot\frac{1}{\sqrt{\delta}}.
\]
In particular, for fixed $\delta$, each $\varepsilon_\bullet = O(S^{-1/2})$.
\end{corollary}

\begin{proof}
From Lemma~\ref{lem:semantic_cancellation}, with probability at least $1-\delta$,
\[
\|R_{\mathrm{sem}}\mathbf{1}\|_2 \le \sqrt{\frac{S}{\delta}}\,\sigma_u.
\]
From Assumption~\ref{assump:agreement_reinforce},
\[
\|A_{\mathrm{agr}}\mathbf{1}\|_2 \ge S\,m_u.
\]
Divide the first inequality by the second to get
\[
\frac{\|R_{\mathrm{sem}}\mathbf{1}\|_2}{\|A_{\mathrm{agr}}\mathbf{1}\|_2}
\le
\frac{\sqrt{S/\delta}\,\sigma_u}{S\,m_u}
=
\frac{1}{\sqrt{S}}\cdot\frac{\sigma_u}{m_u}\cdot\frac{1}{\sqrt{\delta}}
=
\varepsilon_u.
\]
The other two inequalities are identical.
\end{proof}

\subsection{Agreement spike recovery via the averaging estimator (full proof)}
\label{sec:agreement_recovery}

Define
\begin{equation}
\label{eq:uvs_def}
u \triangleq G\mathbf{1},
\qquad
v \triangleq G^\top \mathbf{1},
\qquad
s \triangleq \mathbf{1}^\top G \mathbf{1},
\end{equation}
and the rank-one estimator
\begin{equation}
\label{eq:Ghat1_def}
\widehat{G}_1 \triangleq \frac{u v^\top}{s}.
\end{equation}

\begin{theorem}[Recovery of the rank-one agreement component]
\label{thm:rank_one_agreement}
Assume \eqref{eq:rank_one_decomp}--\eqref{eq:rank_one_A} and $\alpha\neq 0$, $\beta\neq 0$ in
\eqref{eq:alpha_beta}. Suppose \eqref{eq:residual_cancel} holds for some $0\le\varepsilon<1$
(where for simplicity we take $\varepsilon\triangleq\max\{\varepsilon_u,\varepsilon_v,\varepsilon_s\}$).
Define
\begin{equation}
\label{eq:c_def}
c \triangleq \frac{\sigma^2 \alpha \beta}{s}.
\end{equation}
Then
\begin{align}
\label{eq:rank_one_bound}
\|\widehat{G}_1 - c\,ab^\top\|_2
&\le
\frac{2\varepsilon+\varepsilon^2}{1-\varepsilon}\,\|A_{\mathrm{agr}}\|_2,
\\[3pt]
\label{eq:scale_bound}
\left|\frac{c}{\sigma}-1\right|
&\le
\frac{\varepsilon}{1-\varepsilon}.
\end{align}
\end{theorem}

\begin{proof}
Define perturbations
\[
r_u \triangleq R_{\mathrm{sem}}\mathbf{1},\qquad
r_v \triangleq R_{\mathrm{sem}}^\top\mathbf{1},\qquad
r_s \triangleq \mathbf{1}^\top R_{\mathrm{sem}}\mathbf{1}.
\]
Then using $G=A_{\mathrm{agr}}+R_{\mathrm{sem}}$ and $A_{\mathrm{agr}}=\sigma ab^\top$,
\begin{align*}
u &= G\mathbf{1} = A_{\mathrm{agr}}\mathbf{1} + r_u
= \sigma a(b^\top\mathbf{1}) + r_u
= \sigma\beta a + r_u,\\
v &= G^\top\mathbf{1} = A_{\mathrm{agr}}^\top\mathbf{1} + r_v
= \sigma b(a^\top\mathbf{1}) + r_v
= \sigma\alpha b + r_v,\\
s &= \mathbf{1}^\top G\mathbf{1}
= \mathbf{1}^\top A_{\mathrm{agr}}\mathbf{1} + r_s
= \sigma\alpha\beta + r_s.
\end{align*}

From \eqref{eq:residual_cancel} and $A_{\mathrm{agr}}\mathbf{1}=\sigma\beta a$,
$A_{\mathrm{agr}}^\top\mathbf{1}=\sigma\alpha b$, $\mathbf{1}^\top A_{\mathrm{agr}}\mathbf{1}=\sigma\alpha\beta$,
we have
\[
\|r_u\|\le \varepsilon \|\sigma\beta a\|,
\qquad
\|r_v\|\le \varepsilon \|\sigma\alpha b\|,
\qquad
|r_s|\le \varepsilon |\sigma\alpha\beta|.
\]
In particular,
\begin{equation}
\label{eq:s_lower}
|s|
=
|\sigma\alpha\beta + r_s|
\ge
|\sigma\alpha\beta| - |r_s|
\ge
(1-\varepsilon)|\sigma\alpha\beta|.
\end{equation}

Now expand $\widehat{G}_1$:
\[
\widehat{G}_1
=
\frac{(\sigma\beta a + r_u)(\sigma\alpha b + r_v)^\top}{s}.
\]
Subtract $c\,ab^\top$ with $c=\sigma^2\alpha\beta/s$ and expand:
\[
\widehat{G}_1 - c\,ab^\top
=
\frac{\sigma\beta\, a r_v^\top}{s}
+
\frac{\sigma\alpha\, r_u b^\top}{s}
+
\frac{r_u r_v^\top}{s}.
\]
Using $\|xy^\top\|_2=\|x\|_2\|y\|_2$,
\begin{align*}
\|\widehat{G}_1 - c\,ab^\top\|_2
&\le
\frac{\sigma|\beta|\,\|a\|\,\|r_v\|}{|s|}
+
\frac{\sigma|\alpha|\,\|b\|\,\|r_u\|}{|s|}
+
\frac{\|r_u\|\,\|r_v\|}{|s|}.
\end{align*}
Bound each term using \eqref{eq:s_lower}:
\[
\frac{\sigma|\beta|\,\|a\|\,\|r_v\|}{|s|}
\le
\frac{\sigma|\beta|\,\|a\|}{(1-\varepsilon)|\sigma\alpha\beta|}
\cdot
\varepsilon\|\sigma\alpha b\|
=
\frac{\varepsilon}{1-\varepsilon}\,\sigma\|a\|\|b\|.
\]
Similarly,
\[
\frac{\sigma|\alpha|\,\|b\|\,\|r_u\|}{|s|}
\le
\frac{\varepsilon}{1-\varepsilon}\,\sigma\|a\|\|b\|.
\]
For the third term,
\[
\frac{\|r_u\|\,\|r_v\|}{|s|}
\le
\frac{\varepsilon\|\sigma\beta a\|\cdot \varepsilon\|\sigma\alpha b\|}{(1-\varepsilon)|\sigma\alpha\beta|}
=
\frac{\varepsilon^2}{1-\varepsilon}\,\sigma\|a\|\|b\|.
\]
Summing yields
\[
\|\widehat{G}_1 - c\,ab^\top\|_2
\le
\frac{2\varepsilon+\varepsilon^2}{1-\varepsilon}\,\sigma\|a\|\|b\|.
\]
Since $\|A_{\mathrm{agr}}\|_2=\|\sigma ab^\top\|_2=\sigma\|a\|\|b\|$, this proves \eqref{eq:rank_one_bound}.

Finally,
\[
\frac{c}{\sigma}
=
\frac{\sigma\alpha\beta}{\sigma\alpha\beta+r_s}
=
\frac{1}{1+r_s/(\sigma\alpha\beta)}.
\]
Let $x=r_s/(\sigma\alpha\beta)$. Then $|x|\le\varepsilon$ and
\[
\left|\frac{c}{\sigma}-1\right|
=
\left|\frac{1}{1+x}-1\right|
=
\left|\frac{-x}{1+x}\right|
\le
\frac{|x|}{1-|x|}
\le
\frac{\varepsilon}{1-\varepsilon},
\]
proving \eqref{eq:scale_bound}.
\end{proof}

\subsection{Why the probe does not capture intra-sequence semantics (formal statement)}
\label{sec:semantic_non_capture}

The previous results show that the probe statistics $(G\mathbf{1},G^\top\mathbf{1},\mathbf{1}^\top G\mathbf{1})$
are dominated by agreement. We now formalize the complementary statement: \emph{semantic structure can
remain large and correlated within each sequence, yet still be invisible to the probe.}
This is exactly the sense in which intra-sequence semantics remain learnable.

\paragraph{A precise distinction: first-order probe vs.\ semantic energy.}
The probe only depends on three \emph{first-order} aggregated quantities:
\[
R_{\mathrm{sem}}\mathbf{1},\quad R_{\mathrm{sem}}^\top\mathbf{1},\quad \mathbf{1}^\top R_{\mathrm{sem}}\mathbf{1}.
\]
In contrast, intra-sequence semantic learnability is governed by the \emph{token-level} residuals
$\{R_{\mathrm{sem}}^{(s,t)}\}$ entering SGD updates, and can be quantified by \emph{second-order}
(or energy) measures, e.g.
\[
E_{\mathrm{sem}}
\triangleq
\sum_{s=1}^S \sum_{t\in\mathcal{T}_s} \bigl\|R_{\mathrm{sem}}^{(s,t)}\bigr\|_F^2,
\qquad
\text{or}
\qquad
\sum_{t\in\mathcal{T}_s} R_{\mathrm{sem}}^{(s,t)} \bigl(R_{\mathrm{sem}}^{(s,t)}\bigr)^\top.
\]
The key point is: \emph{first-order cancellation does not imply small energy}.

\begin{lemma}[Cancellation does not imply vanishing semantic energy]
\label{lem:energy_vs_sum}
Let $Z_1,\dots,Z_S$ be independent random matrices with $\mathbb{E}[Z_s]=0$.
Then:
\begin{equation}
\label{eq:sum_small_energy_large}
\mathbb{E}\Big\|\sum_{s=1}^S Z_s\Big\|_F^2
=
\sum_{s=1}^S \mathbb{E}\|Z_s\|_F^2.
\end{equation}
In particular, even though $\sum_s Z_s$ concentrates around $0$ in many directions,
the total semantic energy $\sum_s \mathbb{E}\|Z_s\|_F^2$ can scale linearly in $S$.
\end{lemma}

\begin{proof}
Expand the Frobenius norm:
\[
\Big\|\sum_{s=1}^S Z_s\Big\|_F^2
=
\left\langle \sum_{s=1}^S Z_s,\ \sum_{t=1}^S Z_t\right\rangle_F
=
\sum_{s=1}^S \|Z_s\|_F^2
+
2\sum_{s<t} \langle Z_s, Z_t\rangle_F.
\]
Take expectations. By independence and $\mathbb{E}[Z_s]=0$,
\[
\mathbb{E}\langle Z_s, Z_t\rangle_F
=
\langle \mathbb{E}[Z_s],\mathbb{E}[Z_t]\rangle_F
=
0
\qquad (s\neq t).
\]
Thus
\[
\mathbb{E}\Big\|\sum_{s=1}^S Z_s\Big\|_F^2
=
\sum_{s=1}^S \mathbb{E}\|Z_s\|_F^2,
\]
proving \eqref{eq:sum_small_energy_large}.
\end{proof}

\paragraph{Consequences for intra-sequence semantics and learnability.}
Lemma~\ref{lem:semantic_cancellation} shows that the probe responses
$R_{\mathrm{sem}}\mathbf{1}$, $R_{\mathrm{sem}}^\top\mathbf{1}$, and $\mathbf{1}^\top R_{\mathrm{sem}}\mathbf{1}$
scale like $O(\sqrt{S})$ (and hence become negligible relative to agreement, which is $\Theta(S)$).
Lemma~\ref{lem:energy_vs_sum} shows that simultaneously the semantic energy contained in the
sequence-level residuals $\{Z_s\}$ can scale as $\Theta(S)$ and need not vanish.

Therefore, semantic structure can remain large and correlated within each sequence (and thus
useful for learning), while still being \emph{invisible to the specific averaging probe} used to
estimate the agreement spike. This is exactly the desired behavior: the probe extracts the
grammar-driven spike while leaving intra-sequence semantic correlations untouched for SGD to learn.

Although the analysis in this section is framed in terms of the gradient matrix G, the core mechanism—removal of a dominant shared component while preserving high-dimensional variability—is representation-agnostic. Appendix A formalizes this mechanism in a canonical high-dimensional model, where each column corresponds to a token and decomposes into a shared component plus token-specific residuals. The appendix proves that removing the shared mean direction preserves almost all variance of the residuals, thereby justifying the semantic-preservation claim made here for gradients.
\section{Averaging-Based Removal of Shared Linguistic Structure}
\label{sec:avg_remove_linguistic}

We formalize and analyze an averaging-based approximation method for removing
shared linguistic structure from high-dimensional token representations.
The goal is to show that simple averaging removes the dominant, anisotropic
linguistic component while preserving semantic variability, thereby providing
a hardware-friendly alternative to SVD-based spike removal.

\subsection{Model Assumptions}

Let the ambient feature space be $\mathbb{R}^D$, decomposed as
\[
\mathbb{R}^D = \mathcal{L} \oplus \mathcal{S},
\]
where:
\begin{itemize}
    \item $\mathcal{L}$ is a low-dimensional \emph{linguistic subspace}
          of dimension $d_\ell \ll D$, capturing shared grammatical and structural
          regularities common across tokens;
    \item $\mathcal{S}$ is a high-dimensional \emph{semantic subspace}
          of dimension $d_s = D - d_\ell$, capturing token-specific meaning.
\end{itemize}

Let $\{x_i\}_{i=1}^N \subset \mathbb{R}^D$ denote token representations in a batch.
Each token is assumed to decompose as
\begin{equation}
\label{eq:token_decomposition}
x_i = \ell_i + s_i,
\end{equation}
where:
\begin{itemize}
    \item $\ell_i \in \mathcal{L}$ is a \emph{linguistic component}.
          While $\ell_i$ may vary across sequences or contexts, all $\ell_i$
          lie in the same low-dimensional subspace $\mathcal{L}$;
    \item $s_i \in \mathcal{S}$ is a \emph{semantic component}.
\end{itemize}

This formulation explicitly allows linguistic structure to exhibit
\emph{controlled variance} across sequences, while remaining confined to a
low-rank subspace.

\paragraph{Assumption 1 (Semantic compressibility).}
Each $s_i$ lies in $\mathcal{S}$ and is \emph{compressible} in the sense that
its energy is concentrated on a subset of coordinates whose effective dimension
is much smaller than $d_s$.
Equivalently, for any fixed energy fraction $\alpha \in (0,1)$, the number of
coordinates required to capture $\alpha\|s_i\|_2^2$ satisfies
\[
k_\alpha(s_i) \ll d_s.
\]
This reflects that each token activates only a limited number of semantic
features, even though the semantic space itself is large.

\paragraph{Assumption 2 (Semantic diversity under batch sampling).}
Tokens in a batch are drawn from different sentences or sufficiently distant
contexts such that their semantic components $\{s_i\}$ can be treated as
approximately independent samples from a common distribution on $\mathcal{S}$
with finite second moment.
Consequently, the empirical semantic mean concentrates:
\begin{equation}
\label{eq:semantic_mean_concentration}
\left\|
\frac{1}{N}\sum_{i=1}^N (s_i - \mathbb{E}s_i)
\right\|
= O(N^{-1/2})
\quad \text{with high probability}.
\end{equation}

Importantly, semantic components need not have zero mean exactly; it suffices
that any batch-level semantic bias is small compared to the shared linguistic
structure and vanishes under aggregation.

These assumptions are consistent with empirical observations in large language
models: linguistic structure is coherent and low-rank, while semantic information
is high-dimensional, diverse, and weakly correlated across tokens.

\subsection{Second-Moment Matrices}

Define the uncentered second-moment matrix
\begin{equation}
\label{eq:M_def}
M \triangleq \frac{1}{N}\sum_{i=1}^N x_i x_i^\top,
\end{equation}
and the centered covariance matrix
\begin{equation}
\label{eq:C_def}
C \triangleq \frac{1}{N}\sum_{i=1}^N (x_i - \mu)(x_i - \mu)^\top,
\qquad
\mu \triangleq \frac{1}{N}\sum_{i=1}^N x_i.
\end{equation}

\subsection{Effect of Averaging}

Using \eqref{eq:token_decomposition}, we compute the sample mean:
\[
\mu
= \frac{1}{N}\sum_{i=1}^N \ell_i
  + \frac{1}{N}\sum_{i=1}^N s_i.
\]
By Assumption~2, the semantic mean term concentrates and is $O(N^{-1/2})$.
We therefore decompose
\[
\mu = \bar{\ell} + O(N^{-1/2}),
\qquad
\bar{\ell} \triangleq \frac{1}{N}\sum_{i=1}^N \ell_i \in \mathcal{L}.
\]

Substituting into \eqref{eq:M_def}, we obtain
\begin{align}
M
&= \frac{1}{N}\sum_{i=1}^N (\ell_i + s_i)(\ell_i + s_i)^\top \\
&= \underbrace{\frac{1}{N}\sum_{i=1}^N \ell_i \ell_i^\top}_{\text{linguistic spike } L}
 + \underbrace{\frac{1}{N}\sum_{i=1}^N s_i s_i^\top}_{\text{semantic covariance } S}
 + \text{cross terms}.
\end{align}
The cross terms vanish up to $O(N^{-1/2})$ by Assumption~2.
Hence,
\begin{equation}
\label{eq:M_decomp}
M = L + S + O(N^{-1/2}),
\qquad
\mathrm{rank}(L) \le d_\ell.
\end{equation}

Similarly, the centered covariance satisfies the exact identity
\begin{equation}
\label{eq:C_M_relation}
C = M - \mu\mu^\top
  = S + O(N^{-1/2}).
\end{equation}

\subsection{Main Theorem}

\begin{theorem}[Averaging removes shared linguistic structure while preserving semantics]
\label{thm:avg_removes_linguistic}
Under the assumptions above:
\begin{enumerate}
    \item The uncentered second-moment matrix $M$ decomposes into a low-rank
          linguistic spike $L$ with $\mathrm{rank}(L)\le d_\ell$ plus a semantic
          covariance $S$, up to a vanishing $O(N^{-1/2})$ term.
    \item The centered covariance matrix $C$ equals the semantic covariance $S$
          up to the same vanishing term.
    \item For any vector $v \in \mathbb{R}^D$ such that $v \perp \mathcal{L}$,
          \[
          Mv = Cv + O(N^{-1/2}).
          \]
          Hence, $M$ and $C$ induce asymptotically identical linear operators on
          $\mathcal{L}^\perp$.
\end{enumerate}
Consequently, subtracting the mean via averaging removes only the shared
linguistic subspace and preserves all semantic directions relevant for learning.
\end{theorem}

\begin{proof}
The decomposition \eqref{eq:M_decomp} follows from bilinearity and the
concentration of semantic averages.
Equation~\eqref{eq:C_M_relation} is an identity.
For $v \perp \mathcal{L}$, we have $v^\top \ell_i = 0$ for all $i$, and therefore
\[
Mv
= (L + S)v + O(N^{-1/2})
= Sv + O(N^{-1/2})
= Cv + O(N^{-1/2}),
\]
which proves the claim.
\end{proof}

\subsection{Implications for Approximation Algorithms}

Theorem~\ref{thm:avg_removes_linguistic} shows that subtracting the mean
(i.e., simple averaging) removes the dominant low-rank linguistic structure
while leaving semantic covariance---the second-order structure that drives
learning and generalization---essentially unchanged.

In contrast, SVD-based methods (e.g., Metis/Matisse) explicitly compute and remove
the top singular subspace associated with $\mathcal{L}$, which is computationally
and hardware expensive.

Averaging provides a principled approximation to this operation with:
\begin{itemize}
    \item $O(ND)$ cost (single pass over tokens),
    \item no matrix decomposition,
    \item preservation of semantic information up to vanishing error.
\end{itemize}

Thus, averaging-based spike removal offers an efficient and theoretically grounded
alternative to SVD-based linguistic spike suppression in high-dimensional token
representations.

\newpage
\newpage

\begin{comment}
\section{Problem Formulation}

We are given a gradient matrix
\[
G \in \mathbb{R}^{d \times d},
\]
and we seek a low-rank approximation that captures the common structure $A$.
Unlike classical matrix approximation problems, our goal is not to minimize
a reconstruction error over all directions.
Instead, we aim to extract a subspace that reflects the dominant shared factors
driving optimization.

In particular, we are interested in approximations that can be computed using
simple operations such as matrix--vector or matrix--matrix products, without
performing a full singular value decomposition.

\section{Averaging as a Structural Probe}


\subsection{Physical Interpretation of Gradient Entries}
\label{sec:physical_meaning_gradient}

We now explain the physical meaning of the gradient matrix entries and clarify how
row and column indices correspond to input and output features in a transformer
linear projection.

\paragraph{Linear projection structure.}
Consider the key projection in a transformer layer,
\[
k = W_k x,
\qquad
W_k \in \mathbb{R}^{d\times d},
\]
where $x \in \mathbb{R}^d$ is the input hidden state and $k \in \mathbb{R}^d$ is the
resulting key vector.
Written component-wise,
\[
k_i = \sum_{\ell=1}^d W_{k,i\ell}\, x_\ell.
\]
Thus, the row index $i$ corresponds to an \emph{output feature} of the key vector,
while the column index $\ell$ corresponds to an \emph{input feature} of the hidden
state.

\paragraph{Meaning of the gradient matrix.}
During training, the gradient of the loss with respect to the projection matrix,
\[
G = \nabla_{W_k} L \in \mathbb{R}^{d\times d},
\]
has the same shape as $W_k$.
Each entry
\[
G_{i\ell} = \frac{\partial L}{\partial W_{k,i\ell}}
\]
measures how sensitively the loss responds to a small change in the weight connecting
input feature $\ell$ to output feature $i$.

Equivalently, $G_{i\ell}$ quantifies the total pressure exerted by the training data
to adjust how input feature $\ell$ contributes to output feature $i$.

\paragraph{Aggregation over tokens.}
In typical training settings, the loss decomposes as a sum over tokens,
\[
L = \sum_{j=1}^n L^{(j)}.
\]
Accordingly, the gradient aggregates contributions across all tokens in the batch:
\[
G
=
\sum_{j=1}^n
\nabla_{W_k} L^{(j)}
=
\sum_{j=1}^n
\delta_k^{(j)} \bigl(x^{(j)}\bigr)^\top,
\]
where $x^{(j)}$ is the input hidden state of token $j$, and $\delta_k^{(j)}$ is the
corresponding backpropagated error signal at the key projection.

Substituting into the entry-wise expression yields
\[
G_{i\ell}
=
\sum_{j=1}^n
\delta^{(j)}_{k,i}\, x^{(j)}_\ell.
\]
Therefore, each entry $G_{i\ell}$ represents the \emph{aggregate effect, across all
tokens}, of modifying the connection from input feature $\ell$ to output feature $i$.

\paragraph{Row and column interpretations.}
This perspective immediately clarifies the roles of rows and columns:
\begin{itemize}
    \item The $i$-th row $G_{i,:}$ captures how the loss responds to changes in all
    input connections feeding into output feature $i$.
    \item The $\ell$-th column $G_{:,\ell}$ captures how changes in input feature
    $\ell$ affect all output features of the projection.
\end{itemize}

\paragraph{Interpretation of column averaging.}
With this interpretation in mind, the column-averaged vector
\[
u = G\mathbf{1},
\qquad
u_i = \sum_{\ell=1}^d G_{i\ell},
\]
aggregates, for each output feature $i$, the total gradient signal across all input
dimensions.
In other words, column averaging collapses feature-specific input effects and reveals
which output features are being consistently pushed by the training data, regardless
of which input dimensions caused that pressure.

This operation intentionally discards fine-grained semantic detail tied to specific
input features, while preserving coherent, shared structure that manifests across
many inputs and tokens.
As such, column averaging serves as a simple but powerful probe of dominant,
grammar-like gradient structure.


\subsection{Column Averaging}

Let $\mathbf{1} \in \mathbb{R}^d$ denote the vector of all ones.
We define the column-sum (or column-average, up to scaling) vector
\[
u = G \mathbf{1}.
\]
The $i$-th entry of $u$ is
\[
u_i = \sum_{\ell=1}^{d} G_{i\ell},
\]
which aggregates all gradient contributions affecting the $i$-th output
dimension of $W_k$ across input dimensions.

Using the outer-product representation of $G$, we obtain
\[
u = \sum_{j=1}^{n} \delta_k^{(j)} \bigl( (x^{(j)})^\top \mathbf{1} \bigr).
\]
Thus, $u$ is a weighted sum of backpropagated error vectors, where the weight
depends on the total activation of each token’s input representation.

\subsection{Row Averaging}

Symmetrically, we define the row-sum vector
\[
v = G^\top \mathbf{1}.
\]
The $\ell$-th entry of $v$ is
\[
v_\ell = \sum_{i=1}^{d} G_{i\ell},
\]
which aggregates gradient contributions associated with the $\ell$-th input
dimension across all output dimensions.

In terms of token-level quantities,
\[
v = \sum_{j=1}^{n} x^{(j)} \bigl( (\delta_k^{(j)})^\top \mathbf{1} \bigr).
\]
Hence, $v$ highlights input features that consistently drive updates to $W_k$.

\subsection{Rank-One Approximation from Averages}

Given $u$ and $v$, a natural rank-one approximation of $G$ is
\[
\widehat{G}_1 = \frac{u v^\top}{\mathbf{1}^\top G \mathbf{1}}.
\]
This matrix has rank one and matches the total sum of entries of $G$.
It captures the most prominent shared structure that survives both row and
column aggregation.

\section{Assumptions Underlying the Averaging Approximation}
\label{sec:averaging_assumptions}

We now formalize the assumptions under which simple averaging operations reveal
a dominant shared structure while suppressing residual components.

\paragraph{Gradient decomposition.}
Let $G \in \mathbb{R}^{d\times T}$ denote the gradient matrix formed by aggregating
gradients over $T$ tokens (or token groups) within a batch.
We assume that $G$ decomposes as
\begin{equation}
G = A + R,
\end{equation}
where:
\begin{itemize}
    \item $A$ is a \emph{shared (grammar-like)} component that is common across tokens
    and sentences;
    \item $R$ is a \emph{residual (semantic / long-tail)} component that varies across
    tokens and sentences.
\end{itemize}

\paragraph{Rank-one common structure.}
We assume that the shared component $A$ is rank one:
\begin{equation}
A = \sigma a b^\top,
\end{equation}
for some nonzero vectors $a \in \mathbb{R}^d$, $b \in \mathbb{R}^T$, and scalar
$\sigma > 0$.
This models the presence of a single dominant, globally aligned gradient direction,
corresponding to ubiquitous structural patterns such as grammar or formatting.

\paragraph{Coherence with averaging.}
Define the scalar projections of the common structure onto the averaging direction:
\begin{equation}
\alpha \;\triangleq\; \mathbf{1}^\top a,
\qquad
\beta \;\triangleq\; \mathbf{1}^\top b.
\end{equation}
We assume
\begin{equation}
\alpha \neq 0,
\qquad
\beta \neq 0.
\end{equation}
This condition ensures that the common structure does not vanish under row or column
averaging, i.e., $A$ is visible to the probe direction $\mathbf{1}$.

\paragraph{Cancellation of the residual under averaging.}
We now state the key assumption that justifies the averaging approximation.

\begin{assumption}[Residual cancellation]
\label{assump:residual_cancellation}
There exists $0 \le \varepsilon < 1$ such that
\begin{equation}
\|R\mathbf{1}\| \le \varepsilon \|A\mathbf{1}\|,
\qquad
\|R^\top \mathbf{1}\| \le \varepsilon \|A^\top \mathbf{1}\|,
\qquad
|\mathbf{1}^\top R \mathbf{1}| \le \varepsilon |\mathbf{1}^\top A \mathbf{1}|.
\end{equation}
That is, the row sums, column sums, and total sum of the residual are small relative
to those of the shared component.
\end{assumption}

\paragraph{Intuitive example (LLM batching).}
Consider a training batch consisting of many independent sentences with diverse
semantics.
Grammar-related gradients (e.g., syntactic patterns or function words) recur across
sentences and therefore align in a consistent direction, accumulating coherently
into the rank-one component $A$.
In contrast, semantic gradients encode sentence-specific meaning; across unrelated
sentences, these gradients point in different directions and lack global alignment.
Although semantic gradients may be large at the token level, their incoherence causes
them to cancel when aggregated, yielding a small residual under averaging as captured
by Assumption~\ref{assump:residual_cancellation}.



\paragraph{Rigorous justification.}
We now provide a simple argument showing that Assumption~\ref{assump:residual_cancellation}
arises naturally from semantic incoherence.

\begin{proposition}[Cancellation by incoherence]
\label{prop:cancellation}
Let
\begin{equation}
R = [r_1, r_2, \dots, r_T],
\end{equation}
where $r_t \in \mathbb{R}^d$ denotes the semantic gradient contribution of token $t$.
Assume:
\begin{enumerate}
    \item $\mathbb{E}[r_t] = 0$ for all $t$;
    \item the vectors $\{r_t\}_{t=1}^T$ are independent;
    \item $\mathbb{E}\|r_t\|^2 \le \sigma^2$ for all $t$.
\end{enumerate}
Then
\begin{equation}
\mathbb{E}[R\mathbf{1}] = 0,
\qquad
\mathbb{E}\|R\mathbf{1}\|^2 \le T\sigma^2,
\end{equation}
and hence $\|R\mathbf{1}\| = O(\sqrt{T})$ with high probability.
\end{proposition}

\begin{proof}
By definition,
\begin{equation}
R\mathbf{1} = \sum_{t=1}^T r_t.
\end{equation}
Linearity of expectation gives
\begin{equation}
\mathbb{E}[R\mathbf{1}] = \sum_{t=1}^T \mathbb{E}[r_t] = 0.
\end{equation}
Moreover,
\begin{align}
\mathbb{E}\|R\mathbf{1}\|^2
&= \mathbb{E}\Big\|\sum_{t=1}^T r_t\Big\|^2 \\
&= \sum_{t=1}^T \mathbb{E}\|r_t\|^2
    + 2\sum_{s<t} \mathbb{E}[r_s^\top r_t].
\end{align}
Independence and zero mean imply $\mathbb{E}[r_s^\top r_t]=0$ for $s\neq t$, yielding
\begin{equation}
\mathbb{E}\|R\mathbf{1}\|^2 \le T\sigma^2.
\end{equation}
\end{proof}

\paragraph{Consequence.}
If the shared component $A$ is coherent across tokens, then
\begin{equation}
\|A\mathbf{1}\| = \Theta(T),
\qquad
\|A^\top \mathbf{1}\| = \Theta(T),
\qquad
|\mathbf{1}^\top A \mathbf{1}| = \Theta(T),
\end{equation}
whereas the residual satisfies
\begin{equation}
\|R\mathbf{1}\|,\; \|R^\top \mathbf{1}\|,\;
|\mathbf{1}^\top R \mathbf{1}| = O(\sqrt{T}).
\end{equation}
Thus the relative contribution of the residual vanishes as $T$ grows, yielding
Assumption~\ref{assump:residual_cancellation} with
$\varepsilon = O(T^{-1/2})$.


\paragraph{Empirical validation on activation matrices.}
We empirically validate the above averaging intuition on a representative activation matrix
$X\in\mathbb{R}^{(\text{batch}\times\text{seq})\times d}$ obtained from a feedforward layer.
We first confirm that $X$ is strongly anisotropic: the singular value spectrum in
Figure~\hyperref[fig:activation_spectrum]{\ref*{fig:activation_spectrum}} exhibits a small number of dominant modes separated from
the bulk.

To connect this anisotropy to averaging, we compare simple mean-based summaries of $X$ to its
leading rank-one SVD component using cosine similarity (after vectorizing matrices).
Figure~\hyperref[fig:mean_vs_svd]{\ref*{fig:mean_vs_svd}(b)} shows that replacing each entry of $X$ by its \emph{column mean}
(averaging over the token dimension) yields a ``column-mean matrix'' that aligns well with
the leading singular component (cosine similarity about 0.56), whereas the analogous
\emph{row-mean} construction in Figure~\hyperref[fig:mean_vs_svd]{\ref*{fig:mean_vs_svd}(a)} has very low similarity.
Moreover, restricting the same aggregation to a single sentence further strengthens the
column-mean signal: Figure~\hyperref[fig:mean_vs_svd]{\ref*{fig:mean_vs_svd}(d)} increases the alignment to about 0.75,
while the single-sentence row-mean baseline in Figure~\hyperref[fig:mean_vs_svd]{\ref*{fig:mean_vs_svd}(c)} remains weak.
Together, these results indicate that averaging across tokens preserves the dominant shared
direction in $X$, whereas averaging across hidden dimensions does not.

\begin{figure}[ht]
    \centering
    \includegraphics[width=0.75\linewidth]{figures/singular_values_log.png}
    \caption{Singular value spectrum (log scale) of a representative activation matrix $X$, illustrating strong anisotropy.}
    \label{fig:activation_spectrum}
\end{figure}

\begin{figure}[ht]
    \centering
    \begin{minipage}{0.49\linewidth}
        \centering
        \includegraphics[width=\linewidth]{figures/row.png}
        \caption*{(a) Full matrix: row mean}
    \end{minipage}
    \hfill
    \begin{minipage}{0.49\linewidth}
        \centering
        \includegraphics[width=\linewidth]{figures/col.png}
        \caption*{(b) Full matrix: column mean}
    \end{minipage}

    \vspace{0.5em}

    \begin{minipage}{0.49\linewidth}
        \centering
        \includegraphics[width=\linewidth]{figures/row_oneseq.png}
        \caption*{(c) Single sentence: row mean}
    \end{minipage}
    \hfill
    \begin{minipage}{0.49\linewidth}
        \centering
        \includegraphics[width=\linewidth]{figures/col_oneseq.png}
        \caption*{(d) Single sentence: column mean}
    \end{minipage}

    \caption{Cosine similarity between mean-based summaries of $X$ and low-rank SVD reconstructions.
    Panels (a,b) use the full matrix; panels (c,d) use a single-sentence submatrix.
    Column means (b,d) align well with the leading singular component (about $0.56$ and $0.75$ respectively), while row means (a,c) exhibit low similarity.}
    \label{fig:mean_vs_svd}
\end{figure}


\section{Theoretical Analysis: Rank-One Case}
\label{sec:rank_one_analysis}

This section shows that the averaging-based estimator recovers the shared
rank-one component under the assumptions introduced in
Section~\ref{sec:averaging_assumptions}.
We first provide an intuitive explanation of the mechanism, followed by a
complete formal proof.

\subsubsection{Intuition: Why averaging recovers the shared structure}

Recall from Section~\ref{sec:averaging_assumptions} that the gradient matrix
admits a decomposition
\[
G = A + R,
\]
where $A = \sigma a b^\top$ is a rank-one shared (grammar-like) component and
$R$ is a residual (semantic / long-tail) component whose row sums, column sums,
and total sum are small relative to those of $A$.

If $G$ were exactly rank one, i.e., $G = A$, then the averaging quantities
\[
u = G\mathbf{1}, \qquad
v = G^\top \mathbf{1}, \qquad
s = \mathbf{1}^\top G \mathbf{1}
\]
would satisfy
\[
u = \sigma \beta a, \qquad
v = \sigma \alpha b, \qquad
s = \sigma \alpha \beta,
\]
and the estimator
\[
\widehat{G}_1 = \frac{u v^\top}{s}
\]
would recover $A$ exactly.

When the residual $R$ is present, the same quantities become
\[
u = \sigma \beta a + r_u, \qquad
v = \sigma \alpha b + r_v, \qquad
s = \sigma \alpha \beta + r_s,
\]
where the perturbations $r_u$, $r_v$, and $r_s$ arise from averaging $R$.
By Assumption~\ref{assump:residual_cancellation}, these perturbations are small
relative to the corresponding signal terms.

Substituting these expressions into $\widehat{G}_1$ produces three error terms:
a signal--residual interaction, a residual--signal interaction, and a
residual--residual interaction.
Each term is suppressed by the smallness of the averaged residuals, while the
denominator $s$ remains bounded away from zero.
Consequently, the recovered matrix $\widehat{G}_1$ remains close to the true
rank-one component $A$, and its scale deviates from $\sigma$ by a controlled
amount proportional to the residual strength.

We now formalize this argument.

\subsubsection{Formal analysis}

\begin{theorem}
\label{thm:rank_one}
Let
\[
u \triangleq G\mathbf{1}, \qquad
v \triangleq G^\top \mathbf{1}, \qquad
s \triangleq \mathbf{1}^\top G \mathbf{1},
\]
and define the averaging-based estimator
\[
\widehat{G}_1 \triangleq \frac{u v^\top}{s}.
\]
Under the assumptions of Section~\ref{sec:averaging_assumptions}, define
\[
\alpha \triangleq \mathbf{1}^\top a, \qquad
\beta \triangleq \mathbf{1}^\top b, \qquad
c \triangleq \frac{\sigma^2 \alpha \beta}{s}.
\]
Then
\begin{align}
\label{eq:rank_one_error}
\|\widehat{G}_1 - c\, a b^\top\|_2
&\le
\frac{2\varepsilon + \varepsilon^2}{1-\varepsilon}\,\|A\|_2,
\\[4pt]
\label{eq:scale_error}
\left|\frac{c}{\sigma} - 1\right|
&\le
\frac{\varepsilon}{1-\varepsilon}.
\end{align}
\end{theorem}

\begin{proof}
Recall that $G = A + R$ and $A = \sigma a b^\top$.

\paragraph{Step 1: Expansion of averaged quantities.}
We compute
\begin{align}
u
&= G\mathbf{1}
= A\mathbf{1} + R\mathbf{1}
= \sigma a (b^\top \mathbf{1}) + r_u
= \sigma \beta a + r_u,
\\
v
&= G^\top \mathbf{1}
= A^\top \mathbf{1} + R^\top \mathbf{1}
= \sigma b (a^\top \mathbf{1}) + r_v
= \sigma \alpha b + r_v,
\\
s
&= \mathbf{1}^\top G \mathbf{1}
= \mathbf{1}^\top A \mathbf{1} + \mathbf{1}^\top R \mathbf{1}
= \sigma \alpha \beta + r_s,
\end{align}
where
\[
r_u \triangleq R\mathbf{1}, \qquad
r_v \triangleq R^\top \mathbf{1}, \qquad
r_s \triangleq \mathbf{1}^\top R \mathbf{1}.
\]

By Assumption~\ref{assump:residual_cancellation},
\[
\|r_u\| \le \varepsilon \|\sigma \beta a\|, \qquad
\|r_v\| \le \varepsilon \|\sigma \alpha b\|, \qquad
|r_s| \le \varepsilon |\sigma \alpha \beta|,
\]
and hence
\[
|s| \ge (1-\varepsilon)|\sigma \alpha \beta|.
\]

\paragraph{Step 2: Expansion of the estimator.}
Substituting the above expressions yields
\[
\widehat{G}_1
=
\frac{(\sigma \beta a + r_u)(\sigma \alpha b + r_v)^\top}{s}.
\]
Subtracting $c\,ab^\top$ and expanding gives
\[
\widehat{G}_1 - c\,ab^\top
=
\frac{\sigma \beta\, a r_v^\top}{s}
+
\frac{\sigma \alpha\, r_u b^\top}{s}
+
\frac{r_u r_v^\top}{s}.
\]

\paragraph{Step 3: Bounding the error.}
Using $\|xy^\top\|_2 = \|x\|_2 \|y\|_2$ and the bounds above, we obtain
\begin{align*}
\|\widehat{G}_1 - c\,ab^\top\|_2
&\le
\frac{\sigma|\beta|\,\|a\|\,\|r_v\|}{|s|}
+
\frac{\sigma|\alpha|\,\|b\|\,\|r_u\|}{|s|}
+
\frac{\|r_u\|\,\|r_v\|}{|s|}
\\
&\le
\frac{2\varepsilon + \varepsilon^2}{1-\varepsilon}\,
\sigma\|a\|\|b\|.
\end{align*}
Since $\|A\|_2 = \sigma\|a\|\|b\|$, this proves
\eqref{eq:rank_one_error}.

\paragraph{Step 4: Bounding the scale error.}
Finally,
\[
\frac{c}{\sigma}
=
\frac{\sigma\alpha\beta}{\sigma\alpha\beta + r_s}
=
\frac{1}{1 + r_s/(\sigma\alpha\beta)}.
\]
Using $|r_s/(\sigma\alpha\beta)| \le \varepsilon$ yields
\[
\left|\frac{c}{\sigma} - 1\right|
\le
\frac{\varepsilon}{1-\varepsilon},
\]
which proves \eqref{eq:scale_error}.
\end{proof}

\end{comment}

\section{Beyond Rank One: Low-Dimensional Grammar Subspaces}
\label{sec:beyond_rank_one}

The rank-one analysis in Section~\ref{sec:rank_one_analysis} captures the simplest
case in which a single shared direction dominates the gradient matrix.
In practice, however, grammatical and structural effects are rarely one-dimensional.
Multiple correlated factors---such as positional patterns, syntactic roles, and
common functional tokens---may jointly contribute to the gradient, forming a
low-dimensional subspace.

We therefore generalize the previous analysis to the case where the shared component
$A$ has rank $k \ll d$.

\paragraph{Low-rank decomposition.}
We assume that
\begin{equation}
G = A + R,
\end{equation}
where the common component admits a rank-$k$ factorization
\begin{equation}
A = U \Sigma V^\top,
\qquad
U \in \mathbb{R}^{d\times k},\;
V \in \mathbb{R}^{T\times k},\;
\Sigma \in \mathbb{R}^{k\times k},
\end{equation}
and $R$ is a residual component capturing semantic or long-tail effects.

As before, the goal is not to reconstruct $G$ in all directions, but to identify the
column space of $A$, which represents the dominant shared structure driving
optimization.

---

\section{Rank-$k$ Averaging-Based Approximation}
\label{sec:rank_k_averaging}

In the rank-one case, a single averaging vector $\mathbf{1}$ suffices to probe the
dominant shared direction.
To recover a $k$-dimensional shared subspace, we generalize averaging by using
multiple independent probes.

\paragraph{Randomized probing.}
Let $\Omega \in \mathbb{R}^{T\times k}$ be a random probe matrix whose columns are
independent random vectors (e.g., Gaussian or Rademacher).
We form the sketch
\begin{equation}
Y \;\triangleq\; G \Omega \;\in\; \mathbb{R}^{d\times k}.
\end{equation}

Each column of $Y$ is a linear combination of the columns of $G$, analogous to the
single-vector average $G\mathbf{1}$ used in the rank-one case.
Intuitively, the random probes play the role of multiple independent averaging
patterns.

\paragraph{Subspace extraction.}
We compute an orthonormal basis $Q \in \mathbb{R}^{d\times k}$ for the column space
of $Y$ via a thin QR decomposition:
\begin{equation}
Y = QR.
\end{equation}
We then define the rank-$k$ approximation
\begin{equation}
\widehat{G}_k \;\triangleq\; QQ^\top G.
\end{equation}

This procedure is the standard \emph{randomized range finder}: $Q$ approximates the
column space of $A$ without explicitly computing a singular value decomposition.

---

\section{Theoretical Analysis: Rank-$k$ Case}
\label{sec:rank_k_analysis}

We now state a simple guarantee showing that the above procedure recovers the shared
grammar subspace under assumptions analogous to those used in the rank-one case.

\paragraph{Assumptions.}
Assume:
\begin{enumerate}
    \item The shared component $A$ has rank $k$ and dominates the residual in operator
    norm, i.e.,
    \[
    \|R\|_2 \le \varepsilon \|A\|_2
    \quad \text{for some } 0 < \varepsilon < 1.
    \]
    \item The columns of $\Omega$ are independent, isotropic random vectors.
\end{enumerate}

\paragraph{Guarantee (informal).}
With high probability over the choice of $\Omega$, the column space of $Q$ satisfies
\[
\operatorname{dist}\bigl(\operatorname{range}(Q),\operatorname{range}(U)\bigr)
= O(\varepsilon),
\]
and consequently
\[
\| \widehat{G}_k - A \|_2 \le O(\varepsilon)\,\|A\|_2.
\]

That is, the projector $QQ^\top$ preserves the dominant shared subspace while
suppressing the residual component.

\paragraph{Interpretation.}
This result formalizes the same mechanism observed in the rank-one case:
\begin{itemize}
    \item random averaging probes align with the coherent subspace $A$;
    \item incoherent residual components contribute only weakly to the sketch;
    \item projecting onto the sketched subspace recovers the grammar structure.
\end{itemize}

---

\section{Relation to Randomized Linear Algebra}
\label{sec:rla_relation}

The rank-$k$ averaging procedure is closely related to classical results in
randomized linear algebra, including randomized SVD and block power methods.
From this perspective, averaging vectors act as \emph{subspace probes}, and the
dominant grammar subspace is identified as the stable range of the gradient matrix.

Unlike full SVD, the randomized approach:
\begin{itemize}
    \item avoids explicit eigenvalue decomposition,
    \item requires only matrix--vector or matrix--matrix products,
    \item naturally aligns with large-scale and distributed training settings.
\end{itemize}

---

\section{Practical Implications for Transformer Optimization}
\label{sec:practical_implications}

For transformer layers such as $W_q$, $W_k$, and $W_v$, the above analysis explains
why gradient dynamics are often governed by a small number of shared directions.
These directions correspond to grammatical or structural regularities that persist
across tokens and batches.

At the same time, semantic gradients occupy a high-dimensional, weakly coherent
subspace that is easily suppressed by averaging and second-moment normalization.
This observation motivates subspace-aware optimization methods that explicitly
attenuate dominant shared modes in order to free learning capacity for long-tail
semantic structure.

---

\section{Summary and Outlook}
\label{sec:summary_outlook}

This chapter developed a unified view of low-rank structure in transformer gradient
matrices.
Starting from simple averaging operations, we showed how coherent grammatical
structure can be isolated while suppressing diverse semantic effects.
We then generalized this idea from rank-one approximations to low-dimensional
subspaces using randomized probing.

These results provide a conceptual and mathematical foundation for spectral and
subspace-based optimization methods in large-scale deep learning, which we explore
in subsequent chapters.

\newpage
\appendix

% Appendix-local commands/environments (safe when this file is \input into main.tex)
\newcommand{\R}{\mathbb{R}}
\newcommand{\E}{\mathbb{E}}
\newcommand{\1}{\mathbf{1}}

\makeatletter
\@ifundefined{c@remark}{%
  \theoremstyle{remark}%
  \newtheorem{remark}{Remark}%
}{}
\makeatother


\chapter{A Canonical High-Dimensional Model for Averaging-Based Spike Removal and Semantic Preservation}
\label{app:canonical_model}

This appendix provides a detailed, self-contained mathematical instantiation of the
assumptions used in Section~1.6.
The goal is to rigorously justify two claims:
\begin{enumerate}
    \item averaging recovers the shared grammar-driven spike; and
    \item removing this spike preserves essentially all semantic variability.
\end{enumerate}
The model is deliberately high-dimensional and distributional, reflecting the fact
that semantic information in language models is spread over a wide feature space and
manifests through variance rather than mean structure.

\section{Matrix-as-Tokens Model}

Let $G \in \R^{n\times n}$ be a matrix whose columns correspond to $n$ tokens, each
represented by an $n$-dimensional vector.
Denote the $j$-th column by $g_j \in \R^n$, so
\[
G = [g_1,\dots,g_n].
\]

\subsection{Grammar and Semantic Feature Decomposition}

We assume that each token representation decomposes additively as
\begin{equation}
\label{eq:col_decomp}
g_j = g_{\mathrm{gram}} + s_j,
\qquad j=1,\dots,n,
\end{equation}
where:
\begin{itemize}
\item $g_{\mathrm{gram}} \in \R^n$ is a deterministic vector capturing grammatical or
      structural features shared across all tokens; and
\item $s_j \in \R^n$ is a random vector capturing token-specific semantic information.
\end{itemize}

Stacking the semantic components gives
\[
S \triangleq [s_1,\dots,s_n] \in \R^{n\times n},
\]
and therefore
\begin{equation}
\label{eq:G_decomp}
G = g_{\mathrm{gram}}\1^\top + S.
\end{equation}

\section{Semantic Feature Distribution: Wide but Sparse}

We now formalize the notion that semantics are high-dimensional but sparse.

\begin{assumption}[Wide sparse semantic feature model]
\label{assump:sparse_semantics}
There exists a semantic feature dictionary of size $n$.
Each token $j$ activates a random subset $\mathcal{S}_j \subset \{1,\dots,n\}$ with
$|\mathcal{S}_j| = k \ll n$.
The semantic vector takes the form
\[
s_j = \sum_{i \in \mathcal{S}_j} \xi_{j,i} e_i,
\]
where:
\begin{itemize}
\item $\{e_i\}$ is the standard basis of $\R^n$;
\item $\xi_{j,i}$ are i.i.d.\ random variables with
      $\E[\xi_{j,i}] = 0$ and $\E[\xi_{j,i}^2] = \tau^2 < \infty$;
\item the supports $\mathcal{S}_j$ are drawn independently across tokens.
\end{itemize}
\end{assumption}

\begin{remark}
This model captures the intuition that each token expresses only a small subset of
possible semantic features (entities, attributes, relations), while the full semantic
space is very large.
\end{remark}

\begin{assumption}[Finite semantic energy]
\label{assump:energy}
Semantic vectors satisfy
\[
\E[s_j] = 0,
\qquad
\E\|s_j\|_2^2 = \sigma^2 = k\tau^2.
\]
\end{assumption}

\section{Step 1: Averaging Recovers the Grammar Direction}

Define the column mean
\begin{equation}
\label{eq:col_mean}
\bar g \triangleq \frac{1}{n}G\1
= g_{\mathrm{gram}} + \frac{1}{n}\sum_{j=1}^n s_j.
\end{equation}

\begin{proposition}[Column averaging recovers grammar]
\label{prop:mean_recovers_grammar}
Under Assumptions~\ref{assump:sparse_semantics}--\ref{assump:energy},
\[
\bar g \xrightarrow[]{a.s.} g_{\mathrm{gram}}
\qquad \text{as } n\to\infty.
\]
\end{proposition}

\begin{proof}
Fix a coordinate $i\in\{1,\dots,n\}$.
The $i$-th coordinate of $\frac{1}{n}\sum_{j=1}^n s_j$ is
\[
\frac{1}{n}\sum_{j=1}^n s_{j,i}.
\]
Each $s_{j,i}$ is nonzero with probability $k/n$, has zero mean, and finite variance.
By independence across $j$ and Kolmogorov’s strong law of large numbers,
\[
\frac{1}{n}\sum_{j=1}^n s_{j,i} \to 0 \quad \text{almost surely}.
\]
Since this holds for every coordinate and $n$ is finite at each stage,
$\frac{1}{n}\sum_j s_j \to 0$ almost surely as a vector.
\end{proof}

\section{Step 2: Spike Removal via Averaging-Based Projection}

Define the normalized average direction
\[
\hat g \triangleq \frac{\bar g}{\|\bar g\|_2},
\]
and the rank-one projector
\begin{equation}
\label{eq:P_def}
P \triangleq \hat g \hat g^\top,
\qquad
\widetilde G \triangleq (I-P)G.
\end{equation}

This operation removes the shared component using only a matrix--vector product and
normalization, without computing any singular vectors.

\section{Step 3: What Is Removed and What Is Preserved}

\subsection{Grammar Spike Removal}

\begin{proposition}
\label{prop:remove_grammar}
$\|(I-P)g_{\mathrm{gram}}\|_2 \to 0$ almost surely as $n\to\infty$.
\end{proposition}

\begin{proof}
By Proposition~\ref{prop:mean_recovers_grammar}, $\hat g \to g_{\mathrm{gram}}/\|g_{\mathrm{gram}}\|$.
Let $u = g_{\mathrm{gram}}/\|g_{\mathrm{gram}}\|$.
Then $P \to uu^\top$, and
\[
(I-P)g_{\mathrm{gram}} \to (I-uu^\top)g_{\mathrm{gram}} = 0.
\]
\end{proof}

\subsection{Semantic Energy Preservation}

We now show that spike removal eliminates only a vanishing fraction of semantic
energy.

\begin{theorem}[Semantic variance preservation under sparse semantics]
\label{thm:semantic_energy_preserved}
Under Assumptions~\ref{assump:sparse_semantics}--\ref{assump:energy},
\[
\E\|(I-P)S\|_F^2
=
\E\|S\|_F^2 - O(k),
\]
and therefore
\[
\frac{\E\|(I-P)S\|_F^2}{\E\|S\|_F^2}
=
1 - O\!\left(\frac{k}{n}\right).
\]
\end{theorem}

\begin{proof}
For any unit vector $u$, the projection identity gives
\[
\|(I-uu^\top)S\|_F^2 = \|S\|_F^2 - \|u^\top S\|_2^2.
\]
We estimate $\E\|u^\top S\|_2^2$.
Write
\[u^\top S = [u^\top s_1,\dots,u^\top s_n],
\qquad
\|u^\top S\|_2^2 = \sum_{j=1}^n (u^\top s_j)^2.
\]
Taking expectation,
\[
\E\|u^\top S\|_2^2
=
\sum_{j=1}^n \E(u^\top s_j)^2
=
n\, \E(u^\top s_1)^2.
\]

Since $s_1$ has exactly $k$ nonzero coordinates, each contributing $\tau^2$ variance,
and $u$ is dense with $\|u\|_2=1$, we have
\[
\E(u^\top s_1)^2
=
\E\sum_{i\in\mathcal{S}_1} u_i^2 \xi_{1,i}^2
=
\tau^2 \E\sum_{i\in\mathcal{S}_1} u_i^2
=
\tau^2 \cdot \frac{k}{n}.
\]
Therefore
\[
\E\|u^\top S\|_2^2 = n \cdot \tau^2 \cdot \frac{k}{n} = k\tau^2 = O(k).
\]

On the other hand,
\[
\E\|S\|_F^2 = \sum_{j=1}^n \E\|s_j\|_2^2 = n\sigma^2 = nk\tau^2.
\]
Combining,
\[
\E\|(I-P)S\|_F^2
=
nk\tau^2 - k\tau^2
=
k\tau^2(n-1),
\]
and dividing by $\E\|S\|_F^2$ yields the claimed ratio.
\end{proof}

\begin{corollary}
Each token loses only an $O(k/n)$ fraction of its semantic energy under spike removal.
\end{corollary}

\section{Interpretation and Connection to Learning}

This appendix rigorously establishes that:
\begin{itemize}
\item averaging recovers the shared grammar-driven direction (first-order structure);
\item spike removal subtracts only a single shared direction;
\item high-dimensional, sparse semantic variance is preserved up to vanishing
      relative error.
\end{itemize}

Since stochastic gradient descent depends on second-order variability of token-level
updates rather than their mean, preserving semantic variance is sufficient to
maintain semantic learnability.
Thus, the averaging-based spike removal isolates dominant grammatical agreement
structure while leaving semantic learning intact.

\end{document}